\begin{table}[H]
\centering
\caption{Impact du ciblage de volatilité sur les métriques du portefeuille. Le Sharpe est structurellement invariant sous un scaling de levier par un facteur constant ; l'overlay ne crée donc pas de valeur \emph{ex nihilo}, il ajuste l'exposition à la volatilité cible du mandat. La ligne non leveragée constitue la \emph{vraie} métrique d'edge statistique ; la ligne leveragée reflète l'expression opérationnelle de cet edge aux contraintes de volatilité du client. Le multiplicateur effectif moyen (rapport des volatilités réalisées) est inférieur au plafond $31\times$ car le levier s'abaisse pendant les périodes de volatilité élevée.}
\label{tab:leverage_impact}
\begin{tabular}{lrrrr}
\toprule
Configuration & CAGR & Vol ann. & Sharpe & Max DD \\
\midrule
Pondéré $62/9/9/20$ \emph{sans levier} & 2.68\,\% & 4.46\,\% & 0.60 & -10.75\,\% \\
Pondéré $62/9/9/20$ \emph{avec} ciblage vol $0.37$ / cap $31\times$ & 36.42\,\% & 37.04\,\% & 0.98 & -55.00\,\% \\
Multiplicateur effectif moyen & \multicolumn{4}{c}{\textbf{8.3$\times$} (rapport volatilité leveragée / non leveragée)} \\
\bottomrule
\end{tabular}
\end{table}
